\chapter{Statistika Deskriptif}\label{ch:g11:stat}

Statistika memampatkan sebuah daftar bilangan menjadi beberapa ringkasan yang
bermakna: di mana datanya berada (\omterm{def:g11:stat:mean}{rata-rata}, \omterm{def:g11:stat:median}{median}), dan seberapa lebar
sebarannya (\omterm{def:g11:stat:quartiles}{kuartil}, \omterm{def:g11:stat:variance}{ragam}, \omterm{def:g11:stat:variance}{simpangan baku}). Bab ini mengerjakan semuanya pada
satu contoh yang dipakai terus. Ringkasan yang sama muncul kembali untuk peubah
acak pada \cref{ch:g11:prob}, dan hukum peluangnya pada \cref{ch:g12:sums}.

Di sepanjang bab ini, data yang kita pakai adalah nilai (dari $20$) sebuah
kelas berisi $12$ siswa:
\[
8,\ 9,\ 10,\ 10,\ 11,\ 12,\ 12,\ 12,\ 13,\ 14,\ 16,\ 17 .
\]

\section{Rata-rata dan median}

\begin{definition}[Rata-rata]\label{def:g11:stat:mean}
\emph{Rata-rata}\index{rata-rata} nilai $x_1, x_2, \dots, x_n$ adalah
\[
\bar x = \frac{x_1 + x_2 + \dots + x_n}{n}.
\]
Jika nilai $x_i$ muncul dengan frekuensi $n_i$ (dengan
$n_1 + \dots + n_k = n$), maka
$\bar x = \frac{n_1 x_1 + \dots + n_k x_k}{n}$.
\end{definition}

\begin{definition}[Median]\label{def:g11:stat:median}
Sebuah \emph{median}\index{median} kumpulan data terurut adalah nilai yang
membelahnya menjadi dua: paling sedikit separuh nilainya $\leq$ median, dan
paling sedikit separuhnya $\geq$ median. Dalam praktiknya, urutkan $n$ nilai
itu; jika $n$ ganjil, mediannya adalah nilai di tengah; jika $n$ genap, ambil
\omterm{prop:g10:coordgeom:midpoint}{titik tengah} dua nilai di tengah.
\end{definition}

\begin{example}\label{ex:g11:stat:meanmedian}
Untuk data yang kita pakai, jumlahnya $144$, sehingga $\bar x = \frac{144}{12} =
12$. Kedua nilai di tengahnya (yang ke-$6$ dan ke-$7$ pada daftar terurut)
sama-sama $12$, sehingga \omterm{def:g11:stat:median}{mediannya} juga $12$. \omterm{def:g11:stat:mean}{Rata-rata} dan \omterm{def:g11:stat:median}{median} tidak harus
berimpit: \omterm{def:g11:stat:median}{median} mengabaikan \emph{seberapa jauh} nilai ekstremnya, \omterm{def:g11:stat:mean}{rata-rata}
tidak. Mengganti nilai tertinggi $17$ dengan $20$ menggeser \omterm{def:g11:stat:mean}{rata-ratanya} menjadi
$12.25$ tetapi meninggalkan \omterm{def:g11:stat:median}{mediannya} di $12$.
\end{example}

\section{Kuartil}

\begin{definition}[Kuartil, jangkauan antarkuartil]\label{def:g11:stat:quartiles}
\emph{Kuartil pertama}\index{kuartil} $Q_1$ adalah nilai terkecil yang
membuat paling sedikit seperempat datanya $\leq Q_1$; \emph{kuartil
ketiga} $Q_3$ adalah nilai terkecil yang membuat paling sedikit tiga perempat
datanya $\leq Q_3$. \emph{Jangkauan antarkuartil}\index{jangkauan antarkuartil} adalah $Q_3 - Q_1$: yaitu
lebar paruh tengah datanya.
\end{definition}

\begin{example}
$n = 12$: seperempatnya $3$, sehingga $Q_1$ adalah nilai terurut ke-$3$:
$Q_1 = 10$. Tiga perempatnya $9$, sehingga $Q_3$ adalah nilai ke-$9$:
$Q_3 = 13$. \omterm{def:g11:stat:quartiles}{Jangkauan antarkuartilnya} $3$: jadi paruh tengah kelas itu
berada dalam rentang $3$ angka.
\end{example}

\begin{omfigure}
\begin{tikzpicture}[x=0.62cm]
  \draw[->] (7,0) -- (18.4,0) node[right, font=\small] {nilai};
  \foreach \x in {8,10,12,14,16} {
    \draw (\x,0) -- (\x,-0.12) node[below, font=\small] {\x};
  }
  \fill[omDef] (8,0.25) circle (2pt);
  \fill[omDef] (9,0.25) circle (2pt);
  \fill[omDef] (10,0.25) circle (2pt);
  \fill[omDef] (10,0.55) circle (2pt);
  \fill[omDef] (11,0.25) circle (2pt);
  \fill[omDef] (12,0.25) circle (2pt);
  \fill[omDef] (12,0.55) circle (2pt);
  \fill[omDef] (12,0.85) circle (2pt);
  \fill[omDef] (13,0.25) circle (2pt);
  \fill[omDef] (14,0.25) circle (2pt);
  \fill[omDef] (16,0.25) circle (2pt);
  \fill[omDef] (17,0.25) circle (2pt);
  \draw[omThm, thick] (12,-0.35) -- (12,1.15)
    node[above, font=\small] {$\bar x = 12$};
  \draw[omProp, thick] (10,1.55) -- (10,2.35) -- (13,2.35) -- (13,1.55)
    -- cycle;
  \draw[omProp, thick] (12,1.55) -- (12,2.35);
  \draw[omProp, thick] (8,1.95) -- (10,1.95);
  \draw[omProp, thick] (13,1.95) -- (17,1.95);
  \draw[omProp, thick] (8,1.75) -- (8,2.15);
  \draw[omProp, thick] (17,1.75) -- (17,2.15);
  \node[omProp, font=\small, above] at (10,2.35) {$Q_1$};
  \node[omProp, font=\small, above] at (13,2.35) {$Q_3$};
\end{tikzpicture}

{\small Data yang kita pakai sebagai diagram titik, \omterm{def:g11:stat:mean}{rata-ratanya} (merah), dan
ringkasan lima angka di atasnya (jingga): \omterm{def:g10:functions:extrema}{minimum}, $Q_1$, \omterm{def:g11:stat:median}{median}, $Q_3$,
\omterm{def:g10:functions:extrema}{maksimum}. Kotaknya memuat paruh tengah datanya.}
\end{omfigure}

\section{Ragam dan simpangan baku}

\begin{definition}[Ragam, simpangan baku]\label{def:g11:stat:variance}
\emph{Ragam}\index{ragam} sebuah kumpulan data adalah \omterm{def:g11:stat:mean}{rata-rata}
kuadrat simpangan terhadap \omterm{def:g11:stat:mean}{rata-ratanya}:
\[
V = \frac{1}{n}\sum_{i=1}^{n} (x_i - \bar x)^2,
\]
dan \emph{simpangan baku}\index{simpangan baku} adalah
$\sigma = \sqrt{V}$, yang dinyatakan dalam satuan yang sama dengan datanya.
\end{definition}

\begin{proposition}[Rumus jalan pintas]\label{prop:g11:stat:shortcut}
\omterm{def:g11:stat:variance}{Ragamnya} adalah \omterm{def:g11:stat:mean}{rata-rata} kuadrat dikurangi kuadrat \omterm{def:g11:stat:mean}{rata-rata}:
\[
V = \frac{1}{n}\sum_{i=1}^{n} x_i^2 \;-\; \bar x^{\,2}.
\]
\end{proposition}

\begin{proof}
Jabarkan setiap kuadrat simpangannya:
$(x_i - \bar x)^2 = x_i^2 - 2\bar x\, x_i + \bar x^2$. Setelah dijumlahkan lalu
dibagi $n$:
\[
V = \frac1n \sum x_i^2 - 2\bar x \cdot \underbrace{\frac1n \sum
x_i}_{=\ \bar x} + \frac1n \cdot n\bar x^2
= \frac1n\sum x_i^2 - 2\bar x^2 + \bar x^2
= \frac1n\sum x_i^2 - \bar x^2 .
\]
Identitas yang sama, untuk peubah acak, adalah
\cref{prop:g11:prob:konig}.
\end{proof}

\begin{example}\label{ex:g11:stat:variance}
Untuk data yang kita pakai, simpangan terhadap $\bar x = 12$ adalah $-4$,
$-3$, $-2$, $-2$, $-1$, $0$, $0$, $0$, $1$, $2$, $4$, $5$; kuadratnya
berjumlah
\[
16 + 9 + 4 + 4 + 1 + 0 + 0 + 0 + 1 + 4 + 16 + 25 = 80 .
\]
Karena itu
\[
V = \frac{80}{12} = \frac{20}{3} \approx 6.67,
\qquad
\sigma = \sqrt{20/3} \approx 2.6 \text{ angka}.
\]
Secara kasar, sebuah nilai yang khas berjarak sekitar $2.6$ angka dari
\omterm{def:g11:stat:mean}{rata-ratanya}.
\end{example}

\begin{proposition}[Pergantian satuan]\label{prop:g11:stat:affine}
Jika setiap nilainya diubah menjadi $y_i = a x_i + b$, maka
\[
\bar y = a \bar x + b,
\qquad
V_y = a^2\, V_x,
\qquad
\sigma_y = \abs{a}\,\sigma_x .
\]
\end{proposition}

\begin{proof}
$\bar y = \frac1n \sum (a x_i + b) = a\left(\frac1n\sum x_i\right) +
\frac1n \cdot nb = a\bar x + b$. Lalu
$y_i - \bar y = a(x_i - \bar x)$: setiap simpangannya terskalakan dengan $a$
(pergeseran $b$ saling menghapus), sehingga setiap kuadrat simpangannya
terskalakan dengan $a^2$, dan begitu pula \omterm{def:g11:stat:mean}{rata-ratanya}. Menarik \omterm{def:g11:quad:discriminant}{akarnya} memberi
$\sigma_y = \abs a\,\sigma_x$.
\end{proof}

\begin{example}
Suhu yang tercatat dalam derajat Celsius dengan \omterm{def:g11:stat:mean}{rata-rata} $20$ dan \omterm{def:g11:stat:variance}{simpangan
baku} $3$ dikonversikan ke Fahrenheit lewat $F = 1.8\,C + 32$: \omterm{def:g11:stat:mean}{rata-ratanya}
menjadi $1.8 \times 20 + 32 = 68$, \omterm{def:g11:stat:variance}{simpangan bakunya}
$1.8 \times 3 = 5.4$. Menggeser sebuah kumpulan data tidak mengubah sebarannya;
menskalakannya ikut menskalakan sebarannya.
\end{example}

\begin{method}[Meringkas dan membandingkan kumpulan data]\label{met:g11:stat:compare}
Untuk membandingkan dua kumpulan data, sandingkan: \omterm{def:g11:stat:mean}{rata-rata} (atau \omterm{def:g11:stat:median}{median}) untuk
\emph{kedudukannya}, \omterm{def:g11:stat:variance}{simpangan baku} (atau \omterm{def:g11:stat:quartiles}{jangkauan antarkuartil}) untuk
\emph{sebarannya}. Pasangan (\omterm{def:g11:stat:median}{median}, \omterm{def:g11:stat:quartiles}{jangkauan antarkuartil}) bersifat kekar
terhadap nilai ekstrem; pasangan (\omterm{def:g11:stat:mean}{rata-rata}, \omterm{def:g11:stat:variance}{simpangan baku}) memakai setiap
nilainya dan memasok teori peluang.
\end{method}

\begin{example}
Dua pemanah masing-masing melepas $10$ anak panah; nilainya \omterm{def:g11:stat:mean}{rata-rata} $8.5$
untuk keduanya, tetapi \omterm{def:g11:stat:variance}{simpangan bakunya} $0.5$ dan $2.1$. Tingkatnya sama
secara \omterm{def:g11:stat:mean}{rata-rata} --- tetapi pemanah pertama jauh lebih ajek.
\end{example}

\section{Latihan}

\begin{exercise}[$\star$]\label{exo:g11:stat:1}
Hitunglah \omterm{def:g11:stat:mean}{rata-rata}, \omterm{def:g11:stat:median}{median}, \omterm{def:g11:stat:quartiles}{kuartil} dan \omterm{def:g11:stat:quartiles}{jangkauan antarkuartil} kumpulan
data
\[
5,\ 7,\ 7,\ 8,\ 9,\ 10,\ 11,\ 14 .
\]
\end{exercise}

\begin{exercise}[$\star$]\label{exo:g11:stat:2}
Hitunglah \omterm{def:g11:stat:variance}{ragam} dan \omterm{def:g11:stat:variance}{simpangan baku} kumpulan data
$2,\ 4,\ 4,\ 4,\ 5,\ 5,\ 7,\ 9$ (\omterm{def:g11:stat:mean}{rata-ratanya} dahulu, lalu rumus jalan pintas
\cref{prop:g11:stat:shortcut}).
\end{exercise}

\begin{exercise}[$\star$]\label{exo:g11:stat:3}
Sebuah dadu dilempar $50$ kali; hasilnya diringkas lewat frekuensi:
\begin{center}
\begin{tabular}{c|cccccc}
hasil & 1 & 2 & 3 & 4 & 5 & 6\\
\hline
frekuensi & 6 & 9 & 8 & 10 & 9 & 8
\end{tabular}
\end{center}
Hitunglah \omterm{def:g11:stat:mean}{rata-rata} hasilnya dan \omterm{def:g11:stat:median}{mediannya}.
\end{exercise}

\begin{exercise}[$\star\star$]\label{exo:g11:stat:4}
Sebuah kelas berisi $15$ siswa mempunyai \omterm{def:g11:stat:mean}{rata-rata} nilai $11.2$. Siswa ke-$16$
bergabung dan memperoleh $18$. Berapa \omterm{def:g11:stat:mean}{rata-rata} kelas yang baru?
\end{exercise}

\begin{exercise}[$\star\star$]\label{exo:g11:stat:5}
\omterm{def:g11:stat:mean}{Rata-rata} kumpulan data $6, 8, x, 12, 14$ adalah $10$. Carilah $x$, lalu
hitunglah \omterm{def:g11:stat:variance}{simpangan baku} data lengkapnya.
\end{exercise}

\begin{exercise}[$\star$]\label{exo:g11:stat:6}
Nilai dengan \omterm{def:g11:stat:mean}{rata-rata} $12.4$ dan \omterm{def:g11:stat:variance}{simpangan baku} $2.5$ diskalakan ulang lewat
$y = \frac{x}{2} + 4$. Berikan \omterm{def:g11:stat:mean}{rata-rata} dan \omterm{def:g11:stat:variance}{simpangan baku} yang baru.
\end{exercise}

\begin{exercise}[$\star\star$]\label{exo:g11:stat:7}
Kelompok A ($20$ nilai) berrata-rata $10$; kelompok B ($30$ nilai) berrata-rata
$15$. Hitunglah \omterm{def:g11:stat:mean}{rata-rata} $50$ nilai itu setelah digabungkan. Apakah nilainya
\omterm{prop:g10:coordgeom:midpoint}{titik tengah} $12.5$? Mengapa bukan?
\end{exercise}

\begin{exercise}[$\star\star$]\label{exo:g11:stat:8}
Dua jalur produksi mengisi kantong gula $1$ kg. \omterm{def:g10:proba:sample}{Sampelnya} memberi:
jalur 1: \omterm{def:g11:stat:mean}{rata-rata} $1.002$ kg, \omterm{def:g11:stat:variance}{simpangan baku} $0.006$ kg; jalur 2: \omterm{def:g11:stat:mean}{rata-rata}
$1.001$ kg, \omterm{def:g11:stat:variance}{simpangan baku} $0.019$ kg. Jalur mana yang harus ditera ulang lebih
dahulu, dan mengapa?
\end{exercise}

\begin{exercise}[$\star\star$]\label{exo:g11:stat:9}
Sebuah kumpulan data berisi $n = 10$ nilai mempunyai $\sum x_i = 70$ dan
$\sum x_i^2 = 540$. Hitunglah \omterm{def:g11:stat:mean}{rata-rata}, \omterm{def:g11:stat:variance}{ragam} dan \omterm{def:g11:stat:variance}{simpangan bakunya}.
\end{exercise}

\begin{exercise}[$\star\star\star$]\label{exo:g11:stat:10}
Kumpulan data $1,\ 2,\ 2,\ 3,\ 3,\ 3,\ 4,\ 4,\ 5,\ 23$ memuat sebuah
\emph{pencilan}.
\begin{enumerate}
  \item Hitunglah \omterm{def:g11:stat:mean}{rata-rata} dan \omterm{def:g11:stat:median}{mediannya} dengan dan tanpa nilai $23$.
  \item Hitunglah \omterm{def:g11:stat:quartiles}{jangkauan antarkuartilnya} dengan dan tanpa $23$.
  \item Ringkasan mana yang kekar terhadap pencilannya? Tutup dengan sebuah
        anjuran.
\end{enumerate}
\end{exercise}

\section{Soal: Skor-z, sebuah penggaris semesta}

\begin{problem}[{Soal akhir pekan --- simpangan baku membandingkan
apel dengan jeruk, ketaksamaan Chebyshev mengubah $\sigma$ menjadi
jaminan, dan rata-rata memenangkan medalinya sendiri}]
\label{pb:g11:stat:1}
Apakah $15/20$ pada ujian yang sukar lebih baik daripada $13/20$ pada ujian
yang mudah? Apakah $23$ dalam daftar bilangan kecil sebuah kebetulan atau
sebuah pencilan? Bilangan mentah tidak dapat mengatakannya --- tetapi setelah
dibagi \omterm{def:g11:stat:variance}{simpangan baku} yang tepat, semuanya berbicara dalam satu bahasa,
yaitu \emph{\omterm{pb:g11:stat:1}{skor-z}}. Soal ini menguasai mesin \omterm{def:g11:stat:variance}{ragamnya}
(\cref{def:g11:stat:variance},
\cref{prop:g11:stat:shortcut}), membangun penggaris semestanya, lalu
membuktikan ketaksamaan mengagumkan yang mengubah $\sigma$ dari sebuah bilangan
pemerian menjadi sebuah jaminan.

\medskip
\noindent\textbf{Bagian I --- Mekanika \omterm{def:g11:stat:variance}{ragam}.}
\begin{enumerate}
  \item Untuk kumpulan data $4, 8, 6, 5, 2$: hitunglah \omterm{def:g11:stat:mean}{rata-ratanya}, lalu
        \omterm{def:g11:stat:variance}{ragamnya} langsung dari definisinya, lalu
        $\sigma$.
  \item Hitunglah ulang \omterm{def:g11:stat:variance}{ragamnya} dengan rumus jalan pintas
        (\cref{prop:g11:stat:shortcut}) --- \omterm{def:g11:stat:mean}{rata-rata} kuadrat
        dikurangi kuadrat \omterm{def:g11:stat:mean}{rata-rata} --- lalu periksa kecocokannya.
  \item Dengan frekuensi: nilai $8$, $10$, $14$ diperoleh oleh
        $3$, $5$, $2$ siswa. Hitunglah \omterm{def:g11:stat:mean}{rata-rata}, \omterm{def:g11:stat:variance}{ragam} dan
        $\sigma$-nya (dua angka desimal).
  \item Suhu mempunyai \omterm{def:g11:stat:mean}{rata-rata} $20$ dan $\sigma = 3$ (derajat
        Celsius). Berikan \omterm{def:g11:stat:mean}{rata-rata} dan $\sigma$-nya dalam Fahrenheit
        ($F = 1.8\,C + 32$; \cref{prop:g11:stat:affine}).
  \item Buktikan kedua aturan yang baru saja dipakai: menambahkan konstanta $b$
        pada setiap nilainya menggeser \omterm{def:g11:stat:mean}{rata-ratanya} sebesar $b$ dan membiarkan
        $\sigma$ tidak berubah; mengalikan dengan $a$ mengalikan \omterm{def:g11:stat:mean}{rata-ratanya}
        dengan $a$ dan $\sigma$ dengan $\abs a$.
\end{enumerate}

\medskip
\noindent\textbf{Bagian II --- Penggaris semesta.}
Yang disebut \emph{skor-z}\index{skor-z} sebuah nilai $x$ pada kumpulan data yang
berrata-rata $\bar x$ dan bersimpangan baku $\sigma$ adalah
$z = \dfrac{x - \bar x}{\sigma}$: yaitu banyaknya \omterm{def:g11:stat:variance}{simpangan baku}
yang memisahkan $x$ dari \omterm{def:g11:stat:mean}{rata-ratanya}.
\begin{enumerate}[resume]
  \item Ujian A berrata-rata $12$ dan $\sigma = 2$; ujian B berrata-rata
        $8$ dan $\sigma = 4$. Alice memperoleh $15$ pada A, Bob $13$
        pada B. Hitunglah kedua \omterm{pb:g11:stat:1}{skor-z-nya}: penampilan siapa yang berdiri
        lebih jauh di atas kerumunannya?
  \item Hitunglah kelima \omterm{pb:g11:stat:1}{skor-z} kumpulan data pada pertanyaan 1,
        lalu periksalah bahwa semuanya berrata-rata $0$ dan bersimpangan
        baku $1$. Jelaskan dari pertanyaan 5 mengapa
        pembakuan \emph{selalu} menghasilkan \omterm{def:g11:stat:mean}{rata-rata} $0$ dan
        $\sigma = 1$.
  \item Sebuah tes bakat berrata-rata $100$ dan $\sigma = 15$. Nilai
        mentah berapa yang berskor $z = 2$? Berapa \omterm{pb:g11:stat:1}{skor-z} nilai $76$?
  \item Seorang dokter anak mencatat berat badan bayi
        sebagai ``$z = -2.5$''. Katakan apa artinya, dan mengapa pertumbuhan
        dipantau lewat \omterm{pb:g11:stat:1}{skor-z}, bukan lewat kilogram, seiring anak itu
        bertambah usia.
  \item Pencilan pada \cref{exo:g11:stat:10} (kumpulan data
        $1, 2, 2, 3, 3, 3, 4, 4, 5, 23$): hitunglah \omterm{def:g11:stat:mean}{rata-ratanya},
        $\sigma$, dan \omterm{pb:g11:stat:1}{skor-z} nilai $23$. Ambang peringatan yang
        klasik adalah $\abs z \geq 3$: bagaimana putusannya?
  \item Hitunglah ulang $\sigma$ \emph{tanpa} nilai $23$, lalu
        \omterm{pb:g11:stat:1}{skor-z} yang akan dipunyai $23$ terhadap \omterm{def:g11:stat:mean}{rata-rata} dan $\sigma$
        data yang bersih itu. Apa yang telah dilakukan pencilan itu terhadap
        penggaris yang justru dimaksudkan untuk mengukurnya, dan apa yang
        disarankan \cref{exo:g11:stat:10} sebagai penangkalnya?
\end{enumerate}

\medskip
\noindent\textbf{Bagian III --- Jaminan Chebyshev.}
\begin{enumerate}[resume]
  \item Buktikan ketaksamaan Bienaymé--Chebyshev untuk sebuah
        kumpulan data: jika sebagian $p$ nilainya memenuhi
        $\abs{x - \bar x} \geq k\sigma$, maka, dengan hanya menyimpan
        suku-suku itu pada jumlah \omterm{def:g11:stat:variance}{ragamnya},
        \[
        \sigma^2 \geq p \cdot (k\sigma)^2,
        \qquad\text{sehingga}\qquad
        p \leq \frac{1}{k^2} .
        \]
  \item Apa yang dijamin ketaksamaan itu tentang bagian
        sebarang kumpulan data yang berjarak lebih dari $2\sigma$ dari
        \omterm{def:g11:stat:mean}{rata-ratanya}? Lebih dari $3\sigma$? Periksalah pernyataan $2\sigma$ itu
        pada kumpulan data pertanyaan 1.
  \item Sebuah pabrik menghasilkan baut dengan \omterm{def:g11:stat:mean}{rata-rata} $50$ mm dan
        $\sigma = 0.1$ mm. Apa yang dijamin Chebyshev tentang
        bagian baut yang berada di luar
        $\intcc{49.7}{50.3}$? (Produksi yang sesungguhnya jauh lebih baik
        --- kurva lonceng tahun berikutnya mengubah $\frac19$ yang
        berhati-hati ini menjadi sepersekian persen.)
  \item Dua dana sama-sama memberi imbal hasil $5\,\%$ per tahun secara
        \omterm{def:g11:stat:mean}{rata-rata}, dengan $\sigma = 2\,\%$ untuk dana A dan $\sigma = 12\,\%$
        untuk dana B. Untuk masing-masing dana, apa kata Chebyshev tentang
        bagian tahun yang imbal hasilnya di bawah $-19\,\%$? Apa nama
        keuangan bagi $\sigma$?
\end{enumerate}

\medskip
\noindent\textbf{Bagian IV --- Medali \omterm{def:g11:stat:mean}{rata-rata}, dan sebuah peringatan.}
\begin{enumerate}[resume]
  \item Misalkan $f(c)$ \omterm{def:g11:stat:mean}{rata-rata} kuadrat simpangan
        $(x_i - c)^2$ terhadap pusat sebarang $c$. Buktikan
        identitas
        \[
        f(c) = \sigma^2 + (\bar x - c)^2 ,
        \]
        lalu simpulkan: \omterm{def:g11:stat:mean}{rata-rata} adalah satu-satunya peminimum \omterm{def:g11:stat:mean}{rata-rata}
        kuadrat simpangan --- itulah medalinya, sepadan dengan peminimuman
        simpangan mutlak oleh \omterm{def:g11:stat:median}{median} pada
        \cref{pb:g10:stats:1}.
  \item Periksalah identitas itu pada kumpulan data pertanyaan 1 dengan
        $c = 6$: hitunglah $f(6)$ secara langsung dan lewat rumusnya.
  \item Ringkasan sama, data berbeda: periksalah bahwa
        $P = \{3, 5, 7, 9, 11\}$ dan $Q = \{1, 7, 7, 7, 13\}$
        berbagi \omterm{def:g11:stat:mean}{rata-rata} yang sama tetapi tidak berbagi \omterm{def:g11:stat:variance}{ragam} yang sama; lalu
        susunlah kumpulan data berbentuk
        $\{7 - a, 7, 7, 7, 7 + a\}$ dengan \omterm{def:g11:stat:mean}{rata-rata} dan
        \omterm{def:g11:stat:variance}{ragam} yang persis sama dengan $P$ --- namun berbentuk sama sekali lain.
  \item Buktikan bahwa $\sigma = 0$ memaksa semua nilainya sama. (Setiap suku
        pada jumlah kuadrat yang bernilai nol pastilah apa?)
  \item Penutup --- tangga sang ahli statistik, satu baris untuk setiap
        anak tangganya: ukuran pemusatan (\omterm{def:g11:stat:mean}{rata-rata}, \omterm{def:g11:stat:median}{median}), ukuran penyebaran
        ($\sigma$, \omterm{def:g11:stat:quartiles}{jangkauan antarkuartil}), penggaris semesta ($z$), jaminannya
        (Chebyshev), dan peringatan pada pertanyaan 18: ringkasan memampatkan,
        dan pemampatan selalu kehilangan --- sedangkan kurva lonceng tahun
        berikutnya menanti untuk mengubah \omterm{pb:g11:stat:1}{skor-z} menjadi peluang yang eksak.

\end{enumerate}
\end{problem}
